\documentclass[10pt,a4paper]{amsart}
\usepackage[margin=19mm]{geometry}
\usepackage{amsmath,amssymb,mathtools}
\usepackage[scr=rsfs,cal=euler]{mathalfa}
\usepackage{xfrac}
\usepackage[unicode,hidelinks,bookmarks=false]{hyperref}
\newcommand{\ie}{i.e.\ }
\newcommand{\cf}{cf.\ }
\newcommand{\resp}{respectively\ }
\newcommand{\Subsection}{\subsection}
\providecommand{\nobreakdash}{\nobreak\hbox{-}}
\setlength{\emergencystretch}{2em}
\multlinegap0pt
\usepackage[shortlabels]{enumitem}
\usepackage{amssymb}
\usepackage{float}
\usepackage{tikz}
\usepackage{mathtools}
\newtheorem{lemma}{Lemma}
\newtheorem{theorem}{Theorem}
\newtheorem{proposition}{Proposition}
\usepackage{cleveref}
\crefname{lemma}{Lemma}{Lemmas}
\crefname{theorem}{Theorem}{Theorems}
\crefname{proposition}{Proposition}{Propositions}
\newcommand{\R}{\mathbb R}
\newcommand{\abs}[1]{\left\lvert #1\right\rvert}
\title{\vspace{-0.6in}An AI-Assisted Solution to the Signed BAR Conjecture: Uniqueness in the Harrison--Reiman Class and a Completely-$\mathcal{S}$ Class Obstruction}
\author{Yiping Lu, Youheng Zhu}
\date{Source: arXiv:2607.03639v1 (CC BY 4.0)}
\begin{document}
\maketitle

\noindent\emph{Source and licence.} \vspace{-0.6in}An AI-Assisted Solution to the Signed BAR Conjecture: Uniqueness in the Harrison--Reiman Class and a Completely-$\mathcal{S}$ Class Obstruction. Version arXiv:2607.03639v1 (CC BY 4.0), distributed by arXiv under the Creative Commons Attribution 4.0 International licence, http://creativecommons.org/licenses/by/4.0/, as recorded in the arXiv licence metadata for this exact version and shown on the abstract page https://arxiv.org/abs/2607.03639v1. Full licence notice: \texttt{licenses/arxiv-2607.03639-CC-BY-4.0.txt}.\par\smallskip

\noindent\textit{Editorial scope.} Counted unit: the article's proposition prop:resolvent-directional (source0.tex lines 1014--1033). The article's complete original proof of it is delivered with it (source0.tex lines 1035--1129, one \texttt{proof} environment of 95 lines). No mathematics is written, repaired or re-derived: the statement and the proof are the article's own lines, in the article's own notation, and no retained line is substituted or re-flowed.  Retained uncounted as the source-local context the proof needs: lines 240--242; lines 244--246; lines 531--552; lines 607--652.  Nothing was omitted from any retained span, so the package carries no omission record.  No retained line contains a citation, so the wrapper carries no bibliography and no bibliography entry.  No retained line contains a figure environment, an image-inclusion command or drawing code, so no image is delivered and the package declares no asset dependency.  Editorial formatting only, in addition: an AMS \texttt{amsart} preamble that restates the article's own declarations and macros used by the retained lines, namely the environments \texttt{lemma}, \texttt{theorem}, \texttt{proposition} on separate counters, together with the standard packages those lines need.  The retained environments print in this document as Lemma 1, Theorem 1, Proposition 1 in order of appearance, whereas the article prints the same items with the numbering of its own full document.  The complete native source of the article as distributed in the arXiv e-print for this exact version is bundled unchanged as \texttt{sources/204406/source0.tex}.
\begin{equation}\label{eq:Mmatrix}
 R_{ii}>0,\qquad R_{ij}\le0\ (i\ne j),\qquad R^{-1}\ge0.
\end{equation}

\begin{equation}\label{eq:stability}
 R^{-1}\mu<0
\end{equation}

\begin{lemma}[Principal block projection]\label{lem:Mmatrix}\label{lem:projection}
The normalized matrix has the Harrison--Reiman form
\begin{equation}\label{eq:HR-form}
 \widehat R=I-P^T,
 \qquad P\ge0,
 \qquad \rho(P)<1.
\end{equation}
Every principal submatrix $R_{AA}$ is a nonsingular $M$-matrix and $R_{AA}^{-1}\ge0$. In particular, for every nonempty $A\subset J$, the active directions $\{d_i:i\in A\}$ are linearly independent. For $A\subset J$, define
\begin{equation}\label{eq:projection-formula}
 L_Av=v-R_A R_{AA}^{-1}v_A,
\end{equation}
with $L_\varnothing$ equal to the identity. If $A=I(x)$, then $L_A=\mathsf L_x$ is the (unique) linear map from $\mathbb R^d$ to $H_x$ such that $\mathsf L_xv-v\in\operatorname{span}\{R_i:i\in A\}$. Moreover,
\begin{equation}\label{eq:projection-kills}
 L_A R_i=0,
 \qquad i\in A,
\end{equation}
and
\begin{equation}\label{eq:projection-uniform-constant}
 C_{\mathsf L}:=\max_{A\subset J}\left\|I-R_A R_{AA}^{-1}\pi_A\right\|<\infty,
\end{equation}
where $\pi_Av=v_A$ and the expression for $A=\varnothing$ is the identity.
\end{lemma}

\begin{theorem}[Regularity of SRBM]\label{thm:regularity-SRBM}
Under the standing assumptions \eqref{eq:Mmatrix}--\eqref{eq:stability}, the
following hold.
\begin{enumerate}[label=\textup{(\roman*)}]
\item For each $x\in E$ and each prescribed Brownian motion there is a pathwise
unique SRBM $Z^x$, and $Z^x$ is strong Markov.

\item The semigroup $(P_t)$ maps $C_0(E)$ into itself and is strongly continuous
there.

\item There is a constant $K_\Gamma<\infty$, depending only on the normalized
reflection data, such that synchronous solutions satisfy, for all $x,y\in E$ and
$t\ge0$,
\begin{equation}\label{eq:pathwise-Lipschitz}
        \sup_{0\le s\le t}\abs{Z_s^x-Z_s^y}
        \le K_\Gamma\abs{x-y}
        \qquad\text{almost surely.}
\end{equation}

\item For every $x\in E$ there is an adapted RCLL derivative process
$\mathsf J_t^x\in\operatorname{Lin}(H_x,\R^d)$, $t\ge0$.  For each fixed
$w\in G_x$, on an event of probability one the derivative
\[
        \partial_w Z_t^x
        :=
        \lim_{\varepsilon\downarrow0}
        \frac{Z_t^{x+\varepsilon w}-Z_t^x}{\varepsilon}
\]
exists for every $t\ge0$.  Moreover, for every fixed $t>0$,
\begin{equation}\label{eq:fixed-time-derivative}
        \partial_w Z_t^x
        =
        \mathsf J_t^x[\mathsf L_xw]
        \qquad\text{almost surely.}
\end{equation}

\item For every $x\in E$ and every fixed $u\in H_x$,
\[
        \abs{\mathsf J_t^x[u]}
        \le
        K_\Gamma\abs{u}
        \qquad
        \text{for }dt\otimes\mathbb P\text{-almost every }(t,\omega).
\]
\end{enumerate}
\end{theorem}

\begin{proposition}[Directional factorization of the resolvent]\label{prop:resolvent-directional}
Let \(x\in E\). There is a bounded linear functional
\(\Lambda_x:H_x\to \mathbb R\) such that, for every \(w\in G_x\), the one sided
directional derivative exists and satisfies
\begin{equation}
\partial_w^+ g(x)
:=
\lim_{\varepsilon\downarrow0}
\frac{g(x+\varepsilon w)-g(x)}{\varepsilon}
=
\Lambda_x(\mathsf L_xw).
\end{equation}
The linear functional $\ell_x(v):=\Lambda_x(\mathsf L_xv)$ is bounded by
\begin{equation}
|\ell_x(v)|
\le
\frac{K_\Gamma C_{\mathsf L}\|\nabla h\|_\infty}{\lambda}|v|
\end{equation}
for all $v\in \mathbb R^d$. If \(i\in I(x)\), then $\ell_x(R_i)=0.$ Note that no continuity or measurability of the map \(x\mapsto \ell_x\) is asserted.
\end{proposition}

\begin{proof}
Fix \(x\in E\). For \(u\in H_x\), define $\Lambda_x(u)
:=
\mathbb E\int_0^\infty e^{-\lambda t}
\nabla h(Z_t^x)\cdot \mathsf J_t^x[u]\,dt .$ We first record the measurability and integrability facts needed to define
\(\Lambda_x\). Since \(\mathsf J^x\) is adapted and RCLL with values in the finite dimensional space
\(\operatorname{Lin}(H_x,\mathbb R^d)\), for every fixed \(u\in H_x\) the process $(t,\omega)\mapsto \mathsf J_t^x[u](\omega)$ is progressively measurable, hence \(\mathcal B([0,\infty))\otimes\mathcal F\)-measurable.


By Theorem \ref{thm:regularity-SRBM}, for each fixed \(u\in H_x\), we have 
\(
|\mathsf J_t^x[u]|\le K_\Gamma |u|\) holds for \(dt\otimes \mathbb P\)-almost every \((t,\omega)\). Hence  $\mathbb E\int_0^\infty e^{-\lambda t}
\left|\nabla h(Z_t^x)\cdot \mathsf J_t^x[u]\right|\,dt
\le
\frac{K_\Gamma\|\nabla h\|_\infty}{\lambda}|u|<\infty $, \emph{i.e.} the integral
defining \(\Lambda_x(u)\) is absolutely convergent with
respect to \(e^{-\lambda t}dt\otimes\mathbb P\), and
\[
|\Lambda_x(u)|
\le
\frac{K_\Gamma\|\nabla h\|_\infty}{\lambda}|u|,
\qquad u\in H_x .
\]
Since \(\mathsf J_t^x\in \operatorname{Lin}(H_x,\mathbb R^d)\), linearity of
\(\Lambda_x\) follows from linearity of \(\mathsf J_t^x\) and the preceding bound.
Thus \(\Lambda_x\) is a bounded linear functional on \(H_x\).

Now fix \(w\in G_x\). For all sufficiently small \(\varepsilon>0\),
\(x+\varepsilon w\in E\). By the definition of \(g\),
\begin{equation}
\frac{g(x+\varepsilon w)-g(x)}{\varepsilon}
=
\mathbb E\int_0^\infty e^{-\lambda t}
\frac{h(Z_t^{x+\varepsilon w})-h(Z_t^x)}{\varepsilon}\,dt .
\end{equation}
The synchronous Lipschitz estimate gives the deterministic domination
\begin{equation}
\left|
\frac{h(Z_t^{x+\varepsilon w})-h(Z_t^x)}{\varepsilon}
\right|
\le
K_\Gamma \|\nabla h\|_\infty |w| ,
\end{equation}
uniformly for all sufficiently small \(\varepsilon>0\), all \(t\ge0\), and all
sample paths.

On the probability one event in \cref{thm:regularity-SRBM} corresponding to this fixed pair
\((x,w)\), the pathwise directional derivative
\(
\partial_w Z_t^x
=
\lim_{\varepsilon\downarrow0}
\frac{Z_t^{x+\varepsilon w}-Z_t^x}{\varepsilon}\) exists for every \(t\ge0\). Since \(h\in C_c^\infty(\mathbb R^d)\), the
mean value formula gives
\[
\frac{h(Z_t^{x+\varepsilon w})-h(Z_t^x)}{\varepsilon}
\longrightarrow
\nabla h(Z_t^x)\cdot \partial_w Z_t^x
\]
for \(dt\otimes\mathbb P\)-almost every \((t,\omega)\). Dominated convergence
with respect to \(e^{-\lambda t}\,dt\otimes\mathbb P\) therefore yields
\begin{equation}\label{eq:resolvent-derivative-process}
\partial_w^+ g(x)
=
\mathbb E\int_0^\infty e^{-\lambda t}
\nabla h(Z_t^x)\cdot \partial_w Z_t^x\,dt .
\end{equation}

For every fixed \(t>0\), Theorem \ref{thm:regularity-SRBM} gives $\partial_w Z_t^x = \mathsf J_t^x[\mathsf L_xw]$ almost surely. For fixed \(x\) and \(w\), the map $(t,\omega)\mapsto \partial_w Z_t^x(\omega)$ is \(\mathcal B([0,\infty))\otimes\mathcal F\)-measurable, since it is the
pointwise limit of the continuous-in-\(t\) difference quotients $\varepsilon^{-1}\bigl(Z_t^{x+\varepsilon w}-Z_t^x\bigr).$ The process \(t\mapsto \mathsf J_t^x[\mathsf L_xw]\) is measurable because \(\mathsf J^x\) is RCLL. Hence the fixed-time almost sure identity can be integrated in \(t\), giving
the identity \(dt\otimes\mathbb P\)-almost everywhere. By Fubini, this identity holds for \(dt\otimes\mathbb P\)-almost every
\((t,\omega)\). The value at \(t=0\) is irrelevant for the time integral.
Substituting into the preceding display gives
\begin{equation}
\partial_w^+ g(x)
=
\mathbb E\int_0^\infty e^{-\lambda t}
\nabla h(Z_t^x)\cdot \mathsf J_t^x[\mathsf L_xw]\,dt
=
\Lambda_x(\mathsf L_xw),
\end{equation}
which proves the factorization.

Finally, by Lemma \ref{lem:projection}, we have $|\mathsf L_xv|\le C_{\mathsf L}|v|,$ holds for all $v\in\mathbb R^d.$ Therefore
\[
|\ell_x(v)|
=
|\Lambda_x(\mathsf L_xv)|
\le
\frac{K_\Gamma\|\nabla h\|_\infty}{\lambda}|\mathsf L_xv|
\le
\frac{K_\Gamma C_{\mathsf L}\|\nabla h\|_\infty}{\lambda}|v|.
\]
If \(i\in I(x)\), then Lemma \ref{lem:projection} gives \(\mathsf L_xR_i=0\), and hence $\ell_x(R_i)=\Lambda_x(\mathsf L_xR_i)=0.$ This completes the proof.
\end{proof}

\end{document}
